\documentclass[border=3pt]{standalone}
% 公共导言区 —— 所有 TikZ 图 \input 本文件后使用
\usepackage[UTF8, fontset=mac]{ctex}
\usepackage{tikz}
\usetikzlibrary{arrows.meta, positioning, calc, shapes.geometric, decorations.pathmorphing, backgrounds, fit}
\definecolor{zhusha}{HTML}{9B2D20}
\definecolor{mohei}{HTML}{1A1A1A}
\definecolor{shenhui}{HTML}{555555}
\definecolor{qianhui}{HTML}{F5F0EC}
\definecolor{lan}{HTML}{1F4E79}
\definecolor{qing}{HTML}{2E7D6E}
\tikzset{
  block/.style={draw=shenhui, fill=white, thick, rounded corners=1.5pt,
                font=\small, align=center, inner sep=4pt, minimum height=7mm},
  core/.style={block, fill=lan!12, draw=lan, font=\small\bfseries},
  periph/.style={block, fill=qing!10, draw=qing},
  power/.style={block, fill=zhusha!8, draw=zhusha},
  note/.style={font=\footnotesize, text=shenhui, align=center},
  lab/.style={font=\footnotesize, text=mohei},
  sig/.style={-{Stealth[length=2.2mm]}, thick, color=mohei},
  fbsig/.style={-{Stealth[length=2.2mm]}, thick, color=zhusha, dashed},
  vec/.style={-{Stealth[length=3mm]}, very thick, color=zhusha},
}

\begin{document}
\begin{tikzpicture}[x=1mm, y=1mm, font=\small]

% ===== LED 支路（上）：PA5 — 电阻 — LED — GND =====
\node[block, minimum width=11mm] (pa5) at (-24, 8) {PA5};
\coordinate (w1) at (-12, 8);
\draw[thick] (pa5.east) -- (w1);
% 电阻（矩形）
\draw[thick, fill=white] (-8, 6.2) rectangle (-2, 9.8);
\node[note, above=1mm] at (-5, 10) {限流电阻};
\draw[thick] (w1) -- (-8, 8) (-2, 8) -- (2, 8);
% LED（三角+横线+箭头）
\draw[thick, fill=zhusha!20, draw=zhusha] (2,5.5) -- (2,10.5) -- (7,8) -- cycle;
\draw[thick, zhusha] (7,5.5) -- (7,10.5);
\draw[-{Stealth}, thin, zhusha] (6.5,3.2) -- (4,0.8);
\draw[-{Stealth}, thin, zhusha] (8.5,3.2) -- (6,0.8);
\node[note, above=1mm] at (5, 10.6) {LD2};
\draw[thick] (7,8) -- (13,8);
% 地
\draw[thick] (13,8) -- (12,6) (13,8) -- (14,6) (11.4,6) -- (14.6,6) (12.2,4.9) -- (13.8,4.9) (12.8,3.8) -- (13.2,3.8);
\node[note, below] at (13, 2.6) {GND};

% ===== 按键支路（下）：PC13 — 按钮 — GND，片内上拉 =====
\node[block, minimum width=11mm] (pc13) at (-24, -12) {PC13};
\draw[thick] (pc13.east) -- (-12, -12);
% 片内上拉（虚线支路向上）
\draw[thick, dashed, shenhui] (-12,-12) -- (-12,-2);
\node[note, anchor=south] at (-12, -1.5) {内部上拉（读 1）};
% 按钮触点
\draw[thick] (-12,-12) -- (-6,-12) circle (0.6);
\draw[thick] (0,-12) circle (0.6) -- (6,-12);
\draw[thick] (-5.4,-11) -- (-0.6,-8.5);
\node[note, above=0.5mm] at (-3, -8) {按下接通};
% 地
\draw[thick] (6,-12) -- (5,-14) (6,-12) -- (7,-14) (3.4,-14) -- (6.6,-14) (4.2,-15.1) -- (5.8,-15.1) (4.8,-16.2) -- (5.2,-16.2);
\node[note, below] at (6, -17.4) {GND（按下读 0）};

% ===== 中断路径（右）=====
\node[block, minimum width=17mm, fill=qing!10, draw=qing] (exti) at (22, -12) {EXTI13\\下降沿检测};
\node[block, minimum width=17mm] (nvic) at (48, -12) {NVIC\\中断控制器};
\node[block, minimum width=17mm, fill=zhusha!8, draw=zhusha] (isr) at (48, 4) {EXTI15\_10\_IRQHandler\\（中断服务函数）};
\draw[sig, dashed, zhusha] (6,-12) -- (-0,-12);
\draw[sig, dashed, zhusha] (6,-19) -- ++(4,0) |- node[pos=0.75, below, note] {电平变化} (exti.west);
\draw[sig] (exti.east) -- node[above, note] {触发} (nvic.west);
\draw[sig] (nvic.north) -- (isr.south);

\end{tikzpicture}
\end{document}
